\scanpage{7}

% loc: top left
\diagram{\begin{tikzcd}[ampersand replacement=\&]E\arrow[r,shift left=.7ex,"u"]\& F\arrow[l,shift left=.7ex,"v"]\end{tikzcd}}

% loc: upper example a
a) $E$ catégorie ponctuelle, $v$ l'unique foncteur $F\to E$. Il a un adjoint à g. \inserted{$u$} sss $F$ a un objet initial, et alors $u(P)=\varnothing_F$. Alors $u$ est pl. fidèle.

% loc: upper example b
b) $F$ catégorie ponctuelle, $u$ l'unique foncteur $E\to F$. Alors l'adjoint à droite $v$ existe sss $E$ a un objet final, $v(P)=e_E$. Alors $v$ est pl. fidèle.

% loc: left of main paragraph
\diagram{\begin{tikzcd}[ampersand replacement=\&,row sep=normal]E\arrow[d,bend right=50,"{(\alpha')}"']\\E_1\arrow[u,hook',"{\alpha\ \text{pl. fid}}"']\end{tikzcd}}

% loc: middle opening
On suppose que $\alpha$ admet un adjoint à \uline{gauche} $\alpha'$, d'où couple $\alpha',\alpha$ ; \deleted{qui soit} de sorte que $\alpha'$ est un passage à une catégorie de fractions. Considérons alors un autre foncteur pl. fidèle $E_1\xrightarrow{\beta} F$, et supposons que $u=\beta\alpha'$ :

% loc: small triangle
\diagram{\begin{tikzcd}[ampersand replacement=\&] E\arrow[rr,"u"]\arrow[dr,"\alpha'"']\&\&F\\\&E_1\arrow[ur,"\beta"']\&\end{tikzcd}}

% loc: middle question
ait un adjoint à \deleted{gauche}\inserted{droite} (cela exige de toutes façons, sans savoir d'où provient $\alpha'$, que $\alpha$ ait un adjoint à droite \inserted{$\unsure{\alpha|\gamma}$}, \deleted{\unsure{et ce}} ici le fait que $\alpha$ pl. fidèle équivaut à $\alpha'$ passage à cat. des fractions !) Est-il vrai alors que $\beta$ ait un adjoint à droite ? On peut supposer que \deleted{$\unsure{\operatorname{Ob} F}$} $\operatorname{Ob}F=\operatorname{Ob}E_1\cup\{a\}$, il est donné alors par deux foncteurs $\varphi:E_1\to(\mathrm{Ens})$, $\psi:\deleted{\unsure{(\mathrm{Ens})}}E_1^{\circ}\unsure{\to|{-}}(\mathrm{Ens})$, savoir

% loc: pair of functors
\[\varphi(X)=\operatorname{Hom}(a,\beta X)\qquad\psi(X)=\operatorname{Hom}(\beta X,a)\]

% loc: before product
et un morphisme bifonctoriel

% loc: product
\[\unsure{\varpi|\omega}\quad\deleted{\operatorname{Hom}(}\ \varphi(X)\times\psi(Y)\longrightarrow\operatorname{Hom}(Y,X)\]

% loc: lower argument
Dire que $\beta$ a un adjoint a droite \inserted{(partout défini)} signifie que \deleted{\unsure{pour}} le foncteur contravariant $\psi$ est représentable. Prenant $\varphi(X)=\varnothing$ pour tout $X$, la donnée du morphisme bifonctoriel $\varpi$ devient triviale. Exprimons que $\beta\alpha'$ a un adjoint à droite, i.e. que le foncteur sur $E$

% loc: lower functor
\[X\longmapsto\operatorname{Hom}_{F}(\beta\alpha'(X),a)=\psi(\alpha'(X))\]

% loc: bottom question and reformulation
est représentable. Donc la question devient simplement : Un foncteur $\psi:E_1^{\circ}\to(\mathrm{Ens})$ tel que $\psi\circ\alpha'$ est rep, est-il représentable ? \deleted{\unsure{On peut identifier|Identifions}}\inserted{Identifiant} $\psi$ à un foncteur sur $E^{\circ}$ qui transforme objets de $\Sigma$ en bij, on trouve : Soit \deleted{$\unsure{Z}$} $\xi\in\operatorname{Ob}E$ tel que le \inserted{contr.} foncteur qu'il représente sur $E$, $X\mapsto\operatorname{Hom}_E(X,\xi)$, transforme les $u\in\Sigma$ \deleted{bijections} en bijections. Alors est-il vrai que $\exists\,\xi_1\in E_1$

% archive sheet number: 6
